% ============ Solutions and hints for exercises (Part 6: §16-§18) ============
% 原书 pp.293-294; 原书对 §16 的 Ex.16.4,16.6,16.8,16.9、§17 的 Ex.17.6、§18 的 Ex.18.2,18.7,18.8 未给解答, 编号按原书跳号.

\section*{\S\ref{sec:16}}

\noindent\textbf{16.1.}\quad $\le$ is easy; for $\ge$ consider a system of parameters of $M'$.

\noindent\textbf{16.2.}\quad $\Hom_A(A/\mathfrak{a}, A/\mathfrak{b}) = 0$ by Theorem~\hyperref[thm:16.9]{9}.

\noindent\textbf{16.3.}\quad For $P \in \Ass(A/I)$, $\grade P \ge k$ is clear from $P \supset I$. If $\grade P > k$ then by
Ex.~\ref{ex:16.2-ex}, $I:P = I$, which is a contradiction.

\noindent\textbf{16.5.}\quad Suppose that $(A, m)$ is local; then if $m \in \Ass(A)$ we have $\depth A = 0$, but if
$\htop P > 0$ and $P \notin \Ass(A)$ then $\depth A_P > 0$, so that $M = A$ gives a counter-
example. For example, $A = k[X, Y, Z]_{(X,Y,Z)}/(X, Y, Z)^2 \cap (Z)$ and $P = (x, z)A$
satisfy these conditions.

\noindent\textbf{16.7.}\quad Let $\xi = (x_1, \ldots, x_r)$ be a maximal $M$-sequence in $m$, and set $M' = M/\sum x_iM$,
then there exists $0 \neq \xi' \in M'$ such that $m\xi' = 0$. Thus $n\xi' = 0$, but $n^\nu \subset mB$
so $n^\nu\xi' = 0$ and $n \in \Ass_B(M')$, therefore $\xi$ is also a maximal $M$-sequence in $n$.

\noindent\textbf{16.10.}\quad (i) For $r = 1$ the proof is similar to Theorem~\ref{thm:14.3}. If $r > 1$, applying the case
$r = 1$ gives that $(a_1, \ldots, a_{r-1})$ is prime, and we can then use induction. (ii) For
any $Q \in \Ass(A)$ we have $QA_P \in \Ass(A_P)$; if we had $P \subset Q$ then by (i), $A_Q$
is an integral domain, a contradiction. Hence $P \nsubseteq Q$. Therefore using
Ex.~\ref{ex:16.8-ex}, we see that $P$ can be generated by an $A$-sequence. For a counter-
example let $A = k[x, y, z] = k[X, Y, Z]/(X(1 - YZ))$; $P = (x, y, z) = (y, z) = (y - y^2z, z)$
is a prime ideal of height $2$, but $y - y^2z$ is a zero-divisor in $A$.

\section*{\S\ref{sec:17}}

\noindent\textbf{17.1.}\quad (b) Let $k$ be a field; then $A = k[X, Y]/(XY, Y^2)$ is a one-dimensional
ring which is not CM.

\noindent\textbf{17.2.}\quad $x^3, y^3$ is an $A$-sequence, hence also an $R$-sequence, so that $R$ is CM. The
ring $k[x^4, x^3y, xy^3, y^4]$ is not CM.

\noindent\textbf{17.3.}\quad By localisation we need only consider the case of an integral domain, and
it then follows from Theorem~\ref{thm:11.5}, (i).

\noindent\textbf{17.4.}\quad We can assume that $A$ is a local ring. Since $A/J$ is CM we have $\dim
A/J = \depth A/J = r$, and if we set $k$ for the residue class field of $A$ then
$\Ext^i_A(k, A/J) = 0$ for $i < r$. Using the exact sequence $0 \to
J^\nu/J^{\nu+1} \to A/J^{\nu+1} \to A/J^\nu \to 0$ and the fact that $J^\nu/J^{\nu+1}$ is isomorphic
to a direct sum of a number of copies of $A/J$ we get by induction that
$\Ext^i_A(k, A/J^\nu) = 0$ for $i < r$.

\noindent\textbf{17.5.}\quad (i) Let $x_1, \ldots, x_s$ be a maximal $A$-sequence in $P$, and extend to a maximal $A$-
sequence in $m$, $x_1, \ldots, x_s, y_1, \ldots, y_t$. There exists $Q \in \Ass_A(A/(x_1, \ldots, x_s))$
containing $P$, so that by Theorem~\hyperref[thm:17.2]{2}, $\dim A/P \ge \dim A/Q \ge \depth
A/(x_1, \ldots, x_s) = s$. (ii) $\dim A - \htop P \ge \dim (A/P) \ge \depth A - \depth (P, A) \ge \depth A -
\depth A_P$.

\section*{\S\ref{sec:18}}

\noindent\textbf{18.1.}\quad Using Ex.~\ref{ex:16.1-ex}, we see that $A$ is CM $\iff B$ is CM. Assuming CM, we need
only use condition ($5'$) of Theorem~\hyperref[thm:18.1]{1} as a criterion.

\noindent\textbf{18.3.}\quad Given a prime ideal $P$ of $A[X]$, by localising $A$ at $P \cap A$ and factoring out
by a system of parameters we reduce to proving that if $(A, m, k)$ is a zero-
dimensional Gorenstein local ring, and $P$ a prime ideal of $A[X]$ such that
$P \cap A = m$ then $B = A[X]_P$ is Gorenstein. Then $P$ is generated by $m$ together
with a monic polynomial $f(X)$, and the image of $f$ in $k[X]$ is irreducible.
Since $f$ is $B$-regular, if we set $C = B/(f)$ then the maximal ideal of $C$ is
$mC$, and $C \cong A[X]/(f)$; this is a free $A$-module of finite rank, so that
$\Hom_A(C/mC, C) = \Hom_A(k, A) \otimes_A C$ (by Ex.~\ref{ex:7.7-ex}) $\cong C/mC$. So $C$ is Goren-
stein, therefore $B$ also.

\noindent\textbf{18.4.}\quad $R/(x^3, y^3) = k[x^3, y^3, x^2y, xy^2]/(x^3, y^3) \cong k[U, V]/(U^2, V^2, UV)$. In this
ring $(0)$ is not irreducible, so that $R$ is not Gorenstein.

\noindent\textbf{18.5.}\quad For $0 \neq a \in A$, the ideal $aA$ is $\cong A/I$ with $I \neq A$, so there exists a non-zero
map $\varphi\colon aA \to k$; viewed as a map $aA \to E$, this extends to $A \to E$, so that
$0 \neq \Im \varphi \subset aE$.

\noindent\textbf{18.6.}\quad We can consider $M$ as a submodule of $E = E_A(k)$. By faithfulness,
$A \subset \Hom_A(M, M) \subset \Hom_A(M, E)$. But $0 \to \Hom_A(E/M, E) \to
\Hom_A(E, E) = A \to \Hom_A(M, E) \to 0$ is exact, hence $E/M = 0$.
